\Askhsh{1766}{4}
\begin{ylh}{1}
    \item 3.6. Κριτήρια ισότητας ορθογώνιων τριγώνων
    \item 5.6. Εφαρμογές στα τρίγωνα.
\end{ylh}
Δίνεται τετράγωνο $\mathrm{A}\mathrm{B}\Gamma\Delta$. Έστω $\mathrm{E}$ το συμμετρικό σημείο του $\mathrm{B}$ ως προς το $\Delta$ και $\mathrm{H}$ είναι το μέσο της $\mathrm{A}\Delta$. Η προέκταση της $\Gamma\Delta$ τέμνει την $\mathrm{A}\mathrm{E}$ στο $\Theta$.
Να αποδείξετε ότι:
\begin{alist}
    \item $\Delta\Theta = \frac{\mathrm{A}\mathrm{B}}{2}$. \monades{8}
    \item Τα τρίγωνα $\mathrm{A}\Delta\Theta$ και $\mathrm{H}\Delta\Gamma$ είναι ίσα. \monades{9}
    \item Η $\Gamma\mathrm{H}$ είναι κάθετη στην $\mathrm{A}\mathrm{E}$. \monades{8}
\end{alist}
